\documentclass[12pt]{article}
%^ Start of document
\usepackage{times}  %Makes text TimesNewRoman
\usepackage{setspace} %Allows you to set single or double space 
\usepackage{biblatex} %Imports biblatex package
\usepackage{graphicx} % Required for inserting images
\usepackage{authblk}
\usepackage[colorlinks=true, urlcolor=blue, linkcolor=blue]{hyperref}
\usepackage{subfig}
\usepackage{float}
\usepackage{tabularx}
\usepackage{multirow}
\usepackage[paperheight=11in,
   paperwidth=8.5in,
   top=20mm,
   bottom=20mm,
   left=40mm,
   right=40mm]{geometry}
\usepackage{moresize}
\usepackage{tikz}
\usepackage{xcolor}
\usepackage{physics}
\usepackage{amssymb}
\usepackage{mathrsfs}
\usepackage{amsmath}
\usepackage{ragged2e}
\usepackage{comment}
\usepackage{xcolor}
\addbibresource{ElectroHW3.bib}
\begin{document}
\singlespacing  %Sets single spacing for whole document
\title{Classical Electrodynamics Exam Part 2}

\author[1]{Zachary Tullier}
\affil[1]{Student\\Physics Department\\ University of Houston - Clear Lake \\Houston, Texas\\U.S}

\maketitle


\section{Textbook Question 1.15} \label{section1}
    \underline{Question : }
    \newline Prove Thomson's Theorem: If a number of surfaces are fixed in position and a given total charge is placed on each surface, then the electrostatic energy in the region bounded by the surfaces is an absolute minimum when the charges are placed so that every surface is an equipotential, as happens when they are conductors. 


    %\newpage
    \underline{Question 1.15 Answer: }
    The theorem must hold for any number of surfaces; therefore we will consider $n$ number of surfaces.
    \newline Every surface for the minimum energy state should be an equipoptential $\Phi_i = c$
    \newline If this theorem holds we expect $W' > W$ where $W'$ is the work done by charges that aren't bounded by the surface.
    \begin{equation} \label{eq: }
        W = \frac{\epsilon}{2} \int \abs{\vec{E}}^2 d^3x     
    \end{equation}
    The above work equation is formulated in appendix A (Section \ref{App A}).
    \newline Apply condition ($W' - W > 0$)
    \begin{equation} \label{eq: }
        W'-W = \frac{\epsilon}{2} \int \abs{\vec{E'}}^2 d^3x - \frac{\epsilon}{2} \int \abs{\vec{E}}^2 d^3x
    \end{equation}
    Simplify
    \begin{equation} \label{eq: }
        W'-W = \frac{\epsilon}{2} \int \left(\abs{\vec{E'}}^2 - \abs{\vec{E}}^2 \right) d^3x
    \end{equation}
    Reformulate algebraically
%    \newline \textcolor{red}{Do inbetween}
    \begin{equation}
        W'-W = \frac{\epsilon}{2} \int \left((\vec{E'} - \vec{E})^2 - 2 \vec{E} \cdot (\vec{E'} - \vec{E}) \right)d^3x
    \end{equation}
    Substitute Definition of Potential Variable
%    \newline \textcolor{red}{Do inbetween}
    \begin{equation} \label{eq: }
        W'-W = \frac{\epsilon}{2} \int \left((\vec{E'} - \vec{E})^2 - 2 \nabla \Phi \cdot (\vec{E'} - \vec{E}) \right)d^3x        
    \end{equation}
    Simplify
    \begin{equation} \label{eq: }
        W'-W = \frac{\epsilon}{2} \int \left((\vec{E'} - \vec{E})^2 - 2 \nabla (\Phi (\vec{E'} - \vec{E}) - \Phi \nabla (\vec{E'} - \vec{E}) \right)d^3x        
    \end{equation}
    Since $\nabla \cdot \vec{E} = \nabla \cdot \vec{E'}$ and $\int_s \vec{E} \cdot n da = \int_s \vec{E'} \cdot n da$ this equation can be simplified
    \begin{equation} \label{eq: }
        W'-W = \frac{\epsilon}{2} \int \left((\vec{E'} - \vec{E})^2  \right)d^3x        
    \end{equation}    
    The value within the integral will always be positive therefore $W' - W$ will always be positive; thus proving the Validity of Thomson's Theorem

    
\newpage
\section{Textbook Question 2.15} \label{section2}
    \underline{Question: }
    \begin{enumerate}
        \item [a)] Shown that the Green function $G(x, y:x', y')$ appropriate for Dirichlet boundary conditions for a square two-dimensional region, $0 \leq x \leq 1$, $0 \leq y \leq 1$, has an expansion 
            \begin{gather*}
                G(x,y;x', y') = 2 \sum_{n=1}^\infty g_n(y,y') sin(n \pi x) sin(n \pi x')
            \end{gather*}
            Where $g_n(y,y')$ satisfies
            \begin{gather*}
                \left(\frac{\partial^2}{\partial{y'}^2} - n^2 \pi^2 \right) g_n(y,y') = -4 \pi \delta(y'-y)
            \end{gather*}
            and 
            \begin{gather*}
                g_n(y,0) = g_n(y,1) = 0
            \end{gather*}   
        \item [b)] Taking for $g_n(y,y')$ appropriate linear combinations of $sin(n \pi y')$ and $cos(n \pi y')$ in the two regions, $y' < y$ and $y' > y$, in accord with the boundary conditions and the discontinuity in slope required by the souce delta function, show that the explicit form of G is 
            \begin{gather*}
                G(x,y;x', y') = 8 \sum_{n=1}^\infty \frac{1}{n sinh(n \pi)} sin(n \pi x') sinh(n \pi x') sinh (n \pi y_<) sinh(n \pi (1-y_>))
            \end{gather*}
            where $y_<(y_>)$ is the smaller (larger) of $y and y'$
    \end{enumerate}
    
    %\newpage
    \underline{Question 2.15, Section a) Answer: }
    \newline To use a Green's function in a Dirichlet boundary, G must satisfy two conditions.
    \newline First is that G vanishes on the boundary of the region of interest. The suggested expansion satisfies this. If $x'=0$ or $x'=1$ then $sin(n\pi x')=0$. $g(y,y') = 0$ when $y'=0$ or $y'=1$.
    \newline The second condition is that the Greens function satisfies the following equation
    \begin{equation} \label{eq: }
        \nabla^2 G(x,y;x',y') = \delta(x-x') \delta(y-y')
    \end{equation}
    Expand $\nabla$
    \begin{equation} \label{eq: 2.1.1}
        \nabla^2 G(x,y;x',y')= \left(\frac{\partial^2}{\partial x^2} + \frac{\partial^2}{\partial y^2} \right) G(x,y;x',y')
    \end{equation}
    Substitute Green's Function expansion
    \begin{equation} \label{eq: 2.1.2}
        \nabla^2 G(x,y;x',y') = \left(\frac{\partial^2}{\partial x^2} + \frac{\partial^2}{\partial y^2} \right) \left(2 \sum_{n=1}^\infty g_n(y,y') sin(n \pi x) sin(n \pi x') \right)
    \end{equation}
    Rearrange $g_n$
    \begin{equation} \label{eq: 2.1.3}
        \left(\frac{\partial^2}{\partial x^2} g_n(y,y') \right) - g_n(y,y') n^2 \pi^2 = -\delta(y'-y)
    \end{equation}
    $g_n$ has no x component; therefore, it's derivative is zero. Simplify and Isolate $g_n$
    \begin{equation} \label{eq: 2.1.4}
        g_n(y,y') = \frac{\delta(y'-y)}{n^2 \pi^2}
    \end{equation}
    Substitute equation \ref{eq: 2.1.4} into equation \ref{eq: 2.1.2}. 
    \begin{equation} \label{eq: 2.1.5}
        \nabla^2 G(x,y;x',y') = \left(\frac{\partial^2}{\partial x^2} + \frac{\partial^2}{\partial y^2} \right) \left(2 \sum_{n=1}^\infty \frac{\delta(y'-y)}{n^2 \pi^2} sin(n \pi x) sin(n \pi x') \right)
    \end{equation}    
    Perform x double derivation of G. Perform y double derivation of G. Substitute back into the equation
%    \newline \textcolor{red}{Explicitly perform derivations}
    \begin{equation} \label{eq: 2.1.6}
        \nabla^2 G(x,y;x',y') = \delta(y-y') 2 \sum_{n=1}^\infty sin(n \pi x) sin(n \pi x')
    \end{equation}
    Substitute delta function property
%    \newline \textcolor{red}{Explicitly cite property}
    \begin{equation} \label{eq: 2.1.7}
        \nabla^2 G(x,y;x',y') = \delta(y-y') \delta(x-x')
    \end{equation}

    %\newpage
    \underline{Question 2.15, Section b) Answer: }
    \newline The question is essentially asking to explicitly find $g_n$. 
    \newline If $g_n(y,y')$ is the appropriate linear combination $sin(n \pi y')$ and $cos(n \pi y')$ as stated in the question then $g_n$ will can be stated as followsm, if $y<y'$
    \begin{equation} \label{eq: 2.2.1}
        g_n(y,y') = A_1 sinh(n \pi y') + B_1 cosh(n \pi y')  
    \end{equation}
    Apply boundary condition ($g_n(y,y') = 0$ if $y'=1$ to evaluate constants
    \begin{equation} \label{eq: 2.2.2}
        0 = A_1 sinh(n \pi) + B_1 cosh(n \pi)  
    \end{equation}
    Apply algebraic rule ($sinh(X) + cosh(x) = e^{\pm x}$)
%    \newline \textcolor{red}{Find name of rule}
    \begin{equation} \label{eq: 2.2.3}
        0 = (A_1 + B_1) e^{n \pi} + (-A_1 + B_1) e^{-n \pi}
    \end{equation}
    Separate Equations
    \begin{equation} \label{eq: 2.2.4}
        (A_1 + B_1) = -e^{-n \pi}
    \end{equation}
    \begin{equation} \label{eq: 2.2.5}
        (-A_1 + B_1) = e^{-n \pi}
    \end{equation}
    Isolate $B_1$ in equation \ref{eq: 2.2.5}
    \begin{equation} \label{eq: 2.2.6}
        B_1 = e^{-n \pi} + A_1
    \end{equation}
    Substitute into equation \ref{eq: 2.2.4}
    \begin{equation} \label{eq: 2.2.7}
        (A_1 + e^{-n \pi} + A_1) = -e^{-n \pi}
    \end{equation}
    Simplify and Isolate $A_1$
    \begin{equation} \label{eq: 2.2.8}
        2 A_1 = -e^{-n \pi} - e^{-n \pi} = -cosh(n \pi)
    \end{equation}
    Substitute into equation \ref{eq: 2.2.5} and isolate $B_1$
    \begin{equation} \label{eq: 2.2.9}
        B_1 = sinh(n \pi)
    \end{equation}
    Substitute both constants into original $g_n$ equation
    \begin{equation} \label{eq: 2.2.10}
        g_n(y,y') = -cosh(n \pi) sinh(n \pi y') + sinh(n \pi) cosh(n \pi y') = sin(n \pi (1-y'))
    \end{equation}
    Repeat the process if if $y>y'$
    \begin{equation} \label{eq: 2.2.11}
        g_n(y,y') = A_2 sinh(n \pi y') + B_2 cosh(n \pi y')  
    \end{equation}
    Apply boundary condition ($g_n(y,y') = 0$ if  $y'=1$ to evaluate constants
    \begin{equation} \label{eq: 2.2.12}
        g_n = A_2 sinh(n \pi (1-y_>) sin(n \pi y_<)
    \end{equation}
    Differentiate
    \begin{equation} \label{eq: 2.2.13}
        g_n = \frac{-sinh(n \pi (1-y_>) sinh(n \pi y_<)}{n \pi sinh(n \pi)}
    \end{equation}
    Substitute
    \begin{equation} \label{eq: 2.2.14}
        G(x,y;x', y') = 8 \sum_{n=1}^\infty \frac{sin(n \pi x') sinh(n \pi x') sinh (n \pi y_<) sinh(n \pi (1-y_>))}{n sinh(n \pi)} 
    \end{equation}



    
\newpage
\section{Textbook Question 3.13} \label{section3}
    \underline{Question: }
    \newline Solve for the potential in Problem 3.1, using the appropriate Green function obtained in the text, and verify that the answer obtained in this way agrees with the direct solution from the differential equation.

    %\newpage
    \underline{Question 3.13 Answer: }
    \newline Problem 3.1 asks for the potential between two concentri spheres of radii a and b where the upper hemisphere of the inner sphere and the lower hemisphere of the outer sphere are maintined at potential V.
    We will use Dirichlet Green's Function
    \begin{equation} \label{eq: 3.1}
        G(x,x') = \sum_{l,m} \frac{4 \pi}{2 l + l} \frac{1}{1-\left(\frac{a}{b} \right)}^{2l + l} \left(r_<^l - \frac{a^{2l + l}}{r_<^{l + 1}} \right) \left(\frac{1}{r_>^{l + 1}} - \frac{r_>^l}{b^{2l + l}} \right) Y_l^m(\Omega) Y_l^{m*}(\Omega')
    \end{equation}
    With a being the radius of the inner sphere and b being the radius of the outer sphere. There is no charge for the spheres, so the Greens function for the potential only includes the surface integral and the boundary surfaces are disconnected
    \begin{equation} \label{eq: 3.2}
        \Phi(x) = -\frac{1}{4 \pi} \left( \int_{r'=a} \Phi(x') \frac{\partial G}{\partial n'} a^2 d\Omega' - \int_{r'=b} \Phi(x') \frac{\partial G}{\partial n'} b^2 d\Omega' \right)
    \end{equation}
    Compute derivatives
    \begin{equation} \label{eq: 3.3}
        \left(\frac{\partial G}{\partial n'} \right)_{n'=-r'=-a} = \frac{- 4 \pi}{a^2} \sum_{l,m} \frac{1}{1-\left(\frac{a}{b} \right)^{2l+1}} \left(\left(\frac{a}{r} \right)^{l+1} - \left(\frac{a}{b} \right)^{l+1} \left(\frac{r}{b} \right)^l \right) Y_l^m(\Omega) Y_l^{m*}(\Omega')
    \end{equation}
    \begin{equation} \label{eq: 3.4}
        \left(\frac{\partial G}{\partial n'} \right)_{n'=-r'=-b} = \frac{- 4 \pi}{b^2} \sum_{l,m} \frac{1}{1-\left(\frac{a}{b} \right)^{2l+1}} \left(\left(\frac{a}{b} \right)^l - \left(\frac{a}{b} \right)^l \left(\frac{a}{r} \right)^{l+1} \right) Y_l^m(\Omega) Y_l^{m*}(\Omega')
    \end{equation}
    Substitute equations \ref{eq: 3.3} and \ref{eq: 3.4} into equation \ref{eq: 3.2}
    \newline Reduce spherical harmonic expansion to Legendre polynomial by restricting $m=0$ due to the azimuthal symmetry. 
    \begin{equation} \label{eq: 3.5}
        \begin{split}
        \Phi(x) = \sum_l \left(\int_{-1}^1 V_a(\zeta) P_l(\zeta) d\zeta) \right) \frac{2l + 1}{2 \left(1 - \frac{a}{b}^{2l+1} \right)} \left(\left(\frac{a}{r} \right)^{l+1} - \left(\frac{a}{b} \right)^{l+1} \left(\frac{r}{b} \right)^l \right) P_l(cos(\theta)) 
        \\ + \sum_l \left(\int_{-1}^1 V_b(\zeta) P_l(\zeta) d\zeta) \right) \frac{2l + 1}{2 \left(1 - \frac{a}{b}^{2l+1} \right)} \left(\left(\frac{a}{b} \right)^l - \left(\frac{a}{b} \right)^l \left(\frac{a}{r} \right)^{l+1} \right) P_l(cos(\theta)
        \end{split}
    \end{equation}
    Where $\zeta = cos(\theta')$, $V_a$ and $V_b$ are the potentials for the inner and outer sphere. Since $V_a=V_b=V$ for $\zeta<0$ this expression reduces to
    \begin{equation} \label{eq: 3.6}
        \begin{split}
        \Phi(x) = \frac{V}{2} + V \sum_{j=1}^\infty \frac{(-1)^{j+1} (4j-1)\Gamma(j-\frac{1}{2}}{4 \sqrt{\pi} j! 
        \\ \left(1-\left(\frac{a}{b} \right)^{4j-1} \right)} \left(\left(1+\left(\frac{a}{b} \right)^{2j-1} \right) \left(\frac{a}{r} \right)^{2j} - \left(1 + \left(\frac{a}{b} \right)^{2j} \right) \left(\frac{a}{b} \right)^{2j-1} \right) P_{2j-1}(cos(\theta)
        \end{split}
    \end{equation}
    The above equation agrees with the solution to Problem 3.1

\newpage
\section{Textbook Question 4.7} \label{section4}
    \underline{Question: }
    \newline A localized distribution of charge has a charge density
    \begin{gather*}
        \rho(r) = \frac{1}{64 \pi} r^2 e^-r sin^2(\theta)
    \end{gather*}
    \begin{enumerate}
        \item [a)] Make a multi-pole expansion of the potential due to this charge density and determine all the non-vanishing multi-pole moments. Write down the potential at large distances as a finite expansion in Legendre polynomials. 
        \item [b)] Determine the potential explicitly at any point in space, and show that near the origin, correct to  $r^2$ inclusive, 
        \begin{gather*}
            \Phi(r) \simeq \frac{1}{4 \pi \epsilon_0} \left(\frac{1}{4} - \frac{r^2}{120} P_2(cos(\theta)) \right)
        \end{gather*}
        \item [c)] If there exists at the origin a nucleus with a quadruple moment $Q=10^{-28} m^2$, determine the magnitude of the interaction energy, assuming that the unit of charge in $\rho(r)$ above is the electronic charge and the unit of length is the hydrogen Bohr radius $a_0= \frac{4 \pi \epsilon_0 \hbar^2}{m e^2} = 0.529 \cross 10^{-10} m$. Express your answer as a frequency by dividing by Plank's constant ($\hbar$).
    \end{enumerate}
    The charge density in this problem is that for the $m \pm 1$ states of the $2p$ level in hydrogen, while the quadrupole interaction is of the same order as found in molecules.
    
    %\newpage
    \underline{Question 4.7, Section a) Answer: }
    \newline The general solution in terms of a multipole expansion of a localized distribution in spherical harmonics is:
    \begin{equation} \label{eq: 4.1.1}
        \Phi(r,\theta,\phi) = \frac{1}{4 \pi \epsilon_0} \sum_{l=0}^\infty \sum_{m=-l}^l \frac{4 \pi}{2l+1} q_{lm} \frac{Y_{lm}(\theta, \phi)}{r^{l+1}}
    \end{equation}
    Where
    \begin{equation} \label{eq: 4.1.2}
        q_{lm} = \int Y_{lm}^*(\theta', \phi') {r'}^l \rho(x') dx'
    \end{equation}
    \newline Azimuhal symmetry in the charge distribution
    \begin{equation} \label{eq: 4.1.3}
        \Phi(r,\theta, \phi) = \frac{1}{4 \pi \epsilon_0} \sum_{l=0}^\infty \sqrt{\frac{4 \pi}{2l+1}} q_l \frac{P_l(cos(\theta))}{r^{l+1}}
    \end{equation}
    Where 
    \begin{equation} \label{eq: 4.1.4}
        q_l = 2 \pi \sqrt{\frac{2l+1}{4 \pi}} \int_0^\pi \int_0^\infty P_l(cos(\theta')) {r'}^{l+2} \rho(r', \theta') sin(\theta') dr' d\theta'
    \end{equation}
    Plug in the charge density
    \begin{equation} \label{eq: 4.1.5}
        q_l = \frac{2 \pi}{64 \pi} \sqrt{\frac{2l + 1}{4 \pi}} \int_{-1}^1 P_l(x')(1-{x'}^2) dx' \int_0^\infty {r'}^{l+4} e^{-r'} dr'
    \end{equation}
    Use integral property
    \begin{equation} \label{eq: 4.1.6}
        q_l = \frac{2 \pi}{64 \pi} \sqrt{\frac{2l + 1}{4 \pi}} (l+4)! \int_{-1}^1 P_l(x')(1-{x'}^2 dx'
    \end{equation}
    Use Legendre Polynomials
    \begin{equation} \label{eq: 4.1.7}
        {x'}^2 = \frac{2}{3} P_2(x') + \frac{1}{3} P_0(x')
    \end{equation}
    Substitute
    \begin{equation} \label{eq: 4.1.8}
        q_l = \frac{2 \pi}{64 \pi} \sqrt{\frac{2l + 1}{4 \pi}} (l+4)! \left(-\frac{2}{3} \int_{-1}^1 P_l(x') P_2(x') dx' + \frac{2}{3} \int_{-1}^1 P_l(x') P_0(x') dx' \right)
    \end{equation}
    Due to orthogonality, every $q_l$ is zero except when $l=0$ or $l=2$
    \begin{equation} \label{eq: 4.1.9}
        q_0 = \sqrt{\frac{1}{4\pi}}
    \end{equation}
    \begin{equation} \label{eq: 4.1.10}
        -3 \sqrt{\frac{5}{\pi}}
    \end{equation}
    Solve for solution
    \begin{equation} \label{eq: 4.1.11}
        \Phi(r,\theta,\phi) = \frac{1}{4 \pi \epsilon_0} \left(\sqrt{4 \pi} q_0 \frac{1}{r} + \sqrt{\frac{\pi}{5}} q_2 \frac{3 cos^2(\theta) - 1}{r^3} \right)
    \end{equation}
    \begin{equation} \label{eq: 4.1.12}
        \Phi(r,\theta, \phi) = \frac{1}{4 \pi \epsilon_0} \left(\frac{1}{r} + \frac{3 - 9 cos^2(\theta)}{r^3} \right)
    \end{equation}
    \begin{equation} \label{eq: 4.1.13}
        \Phi(r,\theta, \phi) = \frac{1}{4 \pi \epsilon_0} \left(\frac{P_0(cos(\theta))}{r} -6 \frac{P_2(cos(\theta))}{r^3} \right)
    \end{equation}
    
    %\newpage
    \underline{Question 4.7, Section b) Answer: }
    \newline The multipole expansion is only valid if the observation point is external to a local charge distribution
    \begin{equation} \label{eq: 4.2.1}
        \Phi(r,\theta,\phi) = \frac{1}{4 \pi \epsilon_0} \int \frac{\rho(x')}{\abs{\vec{x} - \vec{x}'}} dx'
    \end{equation}
    Convert into spherical coordinates
    \begin{equation} \label{eq: 4.2.2}
        \Phi(r,\theta,\phi) = \frac{1}{4 \pi \epsilon_0} \frac{1}{64 \pi} \int_0^{2\pi} \int_0^{\pi} \int_0^\infty \frac{{r'}^2 e^{-r'} sin^2(\theta')}{\abs{\vec{x} - \vec{x}'}} {r'}^2 sin(\theta') dr' d\theta' d\phi'
    \end{equation}
    Expand in spherical harmonics
    \begin{equation} \label{eq: 4.2.3}
        \begin{split}
        \Phi(r,\theta,\phi) = \frac{1}{4 \pi \epsilon_0} \frac{1}{64 \pi} \int_0^{2\pi} \int_0^{\pi} \int_0^\infty {r'}^4 e^{-r'} sin^2(\theta') 4 \pi \sum_(l=0)^\infty \sum_{m=-l}^l \frac{1}{2l+1} \frac{r_<^l}{r_>^{l+1}} Y_{lm}^*(\theta',\phi') Y_{lm}(\theta, \phi) sin(\theta') dr' d\theta' d\phi'
        \end{split}
    \end{equation}
    Due to azimuthal symmetry, $m=0$
    \begin{equation} \label{eq: 4.2.4}
        \Phi(r,\theta,\phi) = \frac{1}{4 \pi \epsilon_0} \frac{1}{32 \pi} \sum__{l=0}^\infty P_l(cos(\theta)) \int_{-1}^1 p_l(x') (1-{x'}^2) dx' \int_0^\infty {r'}^4 e^{-r'} \frac{r_<^l}{r_>^{l+1}} dr'
    \end{equation}
    Use poly property
    \begin{equation} \label{eq: 4.2.5}
        \begin{split}
        \Phi(r,\theta, \phi) = \frac{1}{4 \pi \epsilon_0} \frac{1}{32} \sum_{l=0}^\infty P_l(cos(\theta)) 
        \\ \left(-\frac{2}{3} \int_{-1}^1 P_l(x') P_2(x') dx' + \frac{2}{3} P_l(x') P_0(x') dx' \right) \int_0^\infty {r'}^4 e^{-r'} \frac{r_<^l}{r_>^{l+1}} dr'
        \end{split}
    \end{equation}
    Due to orthogonality, all terms drop out except $l=0$ and $l=2$
    \begin{equation} \label{eq: 4.2.6}
        \Phi(r,\theta, \phi) = \frac{1}{4 \pi \epsilon_0} \frac{1}{32} \left(-\frac{4}{15} P_2(cos(\theta')) \int_{0}^\infty {r'}^4 e^{-r'} \frac{r_<^2}{r_>^3} dr' + \frac{4}{3} \int_0^\infty {r'}^4 e^{-r'} \frac{1}{r_>^{l+1}} dr' \right)
    \end{equation}
    If r'<r then 
    \begin{equation} \label{eq: 4.2.7}
        \Phi(r,\theta, \phi) = \frac{1}{4 \pi \epsilon_0} \frac{1}{32} \left(-\frac{4}{15} P_2(cos(\theta')) \frac{1}{3} \int_{0}^r {r'}^6 e^{-r'} dr' + \frac{4}{3} \frac{1}{r} \int_0^r {r'}^4 e^{-r'} dr' \right)
    \end{equation}
    If r'>r then
    \begin{equation} \label{eq: 4.2.8}
        \Phi(r,\theta, \phi) = \frac{1}{4 \pi \epsilon_0} \frac{1}{32} \left(-\frac{4}{15} P_2(cos(\theta')) r^2 \int_{0}^\infty {r'} e^{-r'} dr' + \frac{4}{3} \int_0^\infty {r'}^3 e^{-r'} dr' \right)
    \end{equation}
    Combine both Cases
    \begin{equation} \label{eq: 4.2.9}
        \begin{split}
        \Phi(r,\theta, \phi) = \frac{1}{4 \pi \epsilon_0} \frac{1}{32} \left(-\frac{4}{15} P_2(cos(\theta')) r^2 \int_{0}^\infty {r'} e^{-r'} dr' + \frac{4}{3} \int_0^\infty {r'}^3 e^{-r'} dr' \right) 
        \\ + \frac{1}{4 \pi \epsilon_0} \frac{1}{32} \left(-\frac{4}{15} P_2(cos(\theta')) \frac{1}{3} \int_{0}^r {r'}^6 e^{-r'} dr' + \frac{4}{3} \frac{1}{r} \int_0^r {r'}^4 e^{-r'} dr' \right)
        \end{split}
    \end{equation}
    Evaluate Integrals
    \begin{equation} \label{eq: 4.2.10}
        \begin{split}
        \Phi(r,\theta,\phi) = \frac{1}{4 \pi \epsilon_0} \left(\frac{1}{24} \left(-e^{-r} \left(r^2 + 6r + 18 + \frac{24}{r} \right) + \frac{24}{r} \right) 
        \\ - P_2(cos(\theta)) \frac{1}{120} \left(e^{-r} \left( -5r^2 -30r -120 -\frac{360}{r} - \frac{720}{r^2} - \frac{720}{r^3} \right) \frac{720}{r^3} \right) \right)
        \end{split}
    \end{equation}
    Close to the origin, we can disregard any expansions above $r^3$ and equation $e^{-r}=0$
    \begin{equation} \label{eq: 4.2.11}
        \Phi(r,\theta,\phi) \approx \frac{1}{4 \pi \epsilon_0} \left(\frac{1}{4} - \frac{r^2}{120} P_2(cos(\theta) \right)
    \end{equation}

    %\newpage
    \underline{Question 4.7, Section c) Answer: }
    \newline The work is defined as
    \begin{equation} \label{eq: 4.3.1}
        W = \int \rho \Phi d^3x
    \end{equation}
    Substitute equation \ref{eq: 4.2.11}
    \begin{equation} \label{eq: 4.3.2}
        W = \int \rho \left(\frac{1}{4 \pi \epsilon_0} \left(\frac{1}{4} - \frac{r^2}{120} P_2(cos(\theta) \right) \right) d^3x
    \end{equation}
    Convert to SI Units
    \begin{equation} \label{eq: 4.3.3}
        W = \int \rho \left(\frac{-e}{4 \pi \epsilon_0 a_0} \left(\frac{1}{4} - \frac{{\frac{r}{a_0}}^2}{120} P_2(cos(\theta) \right) \right) d^3x
    \end{equation}    
    Evaluate Integral
    \begin{equation} \label{eq: 4.3.4}
        W = \left(\frac{-e}{4 \pi \epsilon_0 a_0} \left(\frac{1}{4} q - \frac{e}{240 a_0^2} Q \right) \right)
    \end{equation}
    Since $q = 0$
    \begin{equation} \label{eq: 4.3.5}
        W = \frac{e^2 Q}{480 \pi \epsilon_0 a_0^3} 
    \end{equation}
    To introduce $\hbar$ we introduce the zero-energy fine structure constant $\alpha = \frac{e^2}{4 \pi \epsilon_0 \hbar c}$
    \begin{equation} \label{eq: 4.3.6}
        W = \frac{\alpha c Q}{480 \pi a_0^3 \hbar}
    \end{equation}
    Rearrange and substitute constants to get
    \begin{equation} \label{eq: 4.3.7}
        \frac{W}{\hbar} \approx 1 [MHz]
    \end{equation}

\section{Appendix A: Energy Density formulation without explicit charge reference} \label{App A}
    If a point charge ($q_i$) is brought from infinity to a point ($x_i$) in a region of localized electric fields described by the scalar potential ($\Phi$) the work done (potential energy ($W$)) on the charge is given by
    \begin{equation} \label{eq: }
        W_i = q_i \Phi(x_i)
    \end{equation}
    Using the definition of the potential ($\Phi$) for many charges approaching a singular point
    \begin{equation} \label{eq: }
        \Phi(x) = \frac{1}{4 \pi \epsilon_0} \sum_{j=1}^{n-1} \frac{q_j}{\abs{x-x_j}}
    \end{equation}
    The total potential energy for multiple source charges
    \begin{equation} \label{eq: }
        W = \frac{1}{4 \pi \epsilon_0} \sum_{i=1}^n \sum_{j<i} \frac{q_i q_j}{\abs{x_i-x_j}}
    \end{equation}
    Rearrange for clarity by changing summation bounds with the understanding that any circumstance where $i=j$ is ommitted
    \begin{equation} \label{eq: }
        W = \frac{1}{8 \pi \epsilon_0} \sum_i \sum_j \frac{q_i q_j}{\abs{x_i-x_j}}
    \end{equation}
    Morphing the equation into a continuous charge  distribution by using the Dirac Delta Functions
%    \newline \textcolor{red}{How is this a use of the delta functions}
    \begin{equation} \label{eq: }
        W = \frac{1}{8 \pi \epsilon_0} \int \int \frac{\rho(x) \rho(x')}{\abs{x-x'}} d^3x d^3x'
    \end{equation}
    Substitute the definition of the scalar potential
    \begin{equation} \label{eq: }
        W = \frac{1}{2} \int \rho(x) \Phi(x) d^3x
    \end{equation}
    Use Poissons Equation to get rid of $\rho$.
%    \newline \textcolor{red}{Do inbetween, cite poissons equation}
    \begin{equation} \label{eq:}
        W = \frac{\epsilon_0}{2} \int \Phi \nabla^2 \Phi d^3x
    \end{equation}
    Integrate by parts
%    \newline \textcolor{red}{Do inbetween}
    \begin{equation} \label{eq: }
        W = \frac{\epsilon_0}{2} \int \abs{\nabla \Phi}^2 d^3x = \frac{\epsilon_0}{2} \int \abs{\vec{E}}^2 d^3x = 
    \end{equation}

    In this formulation, we can interpret the energy as being stored in the electric field surrounding the charges. 
    
    
\newpage
\printbibliography
\end{document}
